\documentclass[11pt]{article}
\usepackage[margin=1in]{geometry}
\usepackage[T1]{fontenc}
\usepackage{lmodern}
\usepackage{microtype}
\usepackage{xcolor}
\usepackage{hyperref}
\usepackage{booktabs}
\usepackage{enumitem}
\usepackage{listings}
\usepackage{tcolorbox}
\usepackage{parskip}
\usepackage{array}
\newcolumntype{L}[1]{>{\raggedright\arraybackslash}p{#1}}

% \path{} (from url, loaded by hyperref) prints paths and identifiers
% verbatim -- so `--` is not turned into a dash -- and may break at / _ - .
\makeatletter
\g@addto@macro\UrlBreaks{\do\_\do\-\do\.}
\makeatother
\urlstyle{tt}
\setlength{\emergencystretch}{3em}

\hypersetup{
  colorlinks=true,
  linkcolor=blue!60!black,
  urlcolor=blue!60!black
}

\definecolor{codebg}{RGB}{245,245,245}
\definecolor{codecomment}{RGB}{90,100,90}
\definecolor{codekeyword}{RGB}{0,70,140}
\definecolor{codestring}{RGB}{150,60,20}

\lstdefinestyle{pythonstyle}{
  language=Python,
  backgroundcolor=\color{codebg},
  basicstyle=\ttfamily\small,
  keywordstyle=\color{codekeyword}\bfseries,
  commentstyle=\color{codecomment}\itshape,
  stringstyle=\color{codestring},
  showstringspaces=false,
  breaklines=true,
  frame=single,
  rulecolor=\color{black!20},
  columns=fullflexible,
  keepspaces=true
}
\lstdefinestyle{shellstyle}{
  language=bash,
  backgroundcolor=\color{codebg},
  basicstyle=\ttfamily\small,
  showstringspaces=false,
  breaklines=true,
  frame=single,
  rulecolor=\color{black!20},
  columns=fullflexible,
  keepspaces=true
}
\lstset{style=pythonstyle}

\newtcolorbox{keyidea}{
  colback=blue!5, colframe=blue!50!black,
  title=Key Idea, fonttitle=\bfseries
}
\newtcolorbox{tipbox}{
  colback=green!5, colframe=green!45!black,
  title=Tip, fonttitle=\bfseries
}
\newtcolorbox{warningbox}{
  colback=orange!7, colframe=orange!60!black,
  title=Common Pitfall, fonttitle=\bfseries
}
\title{
  \textbf{Codebase Change Summary --- Week 11}\\[4pt]
  \large System Testing and Acceptance Testing\\[2pt]
  \normalsize SE 413/513 --- Software Testing
}
\author{Demo Testing Suite}
\date{}

\begin{document}
\maketitle

\begin{center}
\textit{What changed in \path{demo-testing-suite/} since Week 10, and why}
\end{center}

\vspace{0.3cm}

% ============================================================
\section{Purpose of This Document}
% ============================================================
The \path{demo-testing-suite/} codebase is a single, shared project that
grows one lecture at a time. This document summarizes what was added or
changed for Week 11 and explains \emph{why} each change exists --- which
concept from the two Week 11 readings, \textit{System Testing: Testing
the Whole System from the Outside} and \textit{Acceptance Testing:
Building the Right Thing}, it demonstrates. Read it alongside the
readings (\path{reading_materials/w11.1_system_testing.pdf},
\path{reading_materials/w11.2_acceptance_testing.pdf}) and
\path{demos/system_acceptance_testing/README.md}.

\begin{keyidea}
Nothing from Weeks 1--10 was removed or changed in behavior. This week
adds one self-contained project for both readings: Week 10's shop,
grown into a system that runs on its own and is tested from the outside
--- as a separate server process, over HTTP only --- at two new levels.
System tests check it against its system requirements; acceptance
tests, written in Gherkin, check it against its user stories. Writing
the system tests found a real defect, present in Week 10's shop too.
\end{keyidea}

% ============================================================
\section{What Changed At a Glance}
% ============================================================
\begin{center}
\small
\begin{tabular}{@{}L{6.0cm}L{1.5cm}L{6.9cm}@{}}
\toprule
\textbf{File} & \textbf{Status} & \textbf{Purpose} \\
\midrule
\multicolumn{3}{@{}l}{\textit{The readings' system (not in the default run)}}\\
\path{demos/system_acceptance_testing/shop/} & New & Week 10's shop
plus \path{GET /health}, \path{GET /products}, coupon codes
(\path{pricing.py}), stricter input validation, and a JSON \texttt{500}
handler. \\
\path{shop/service_racy.py}, \path{shop/pricing_per_unit.py} & New &
Two deliberately defective components: one oversells under concurrency;
one misreads the coupon story. \\
\path{serve.py}, \path{variants.py} & New & Run the shop as one system
(\texttt{--port 0}, \texttt{--db}, \texttt{--seed}); switch defective
components in by environment variable, also inside the server process. \\
\path{fakes/payment_provider.py} & New (copy) & Week 10's fake PayFast,
plus test-only \path{/_test/config} and \path{/_test/charges}. \\
\path{specs/} & New & System requirements, a one-page system test plan,
user stories with acceptance criteria, the API contract v1.1, and the
patch that reproduces finding F1. \\
\path{tests/conftest.py} & New & The system test harness: starts
\path{serve.py} as a subprocess and talks to it over HTTP. \\
\path{tests/system/} & New & 23 system tests: scenarios, hostile input,
access control, concurrency, performance, recovery. \\
\path{tests/acceptance/} & New & 3 \texttt{.feature} files, their step
definitions, and 14 scenarios run by \texttt{pytest-bdd}. \\
\path{tests/unit/}, \path{tests/integration/} & New & The developer's
pricing tests, and Week 10's integration harness carried over: the
levels below, for comparison. \\
\midrule
\multicolumn{3}{@{}l}{\textit{Tooling and documentation}}\\
\texttt{requirements.txt} & Modified & Adds \texttt{pytest-bdd}. \\
\texttt{pytest.ini} & Modified & Registers the \texttt{system} and
\texttt{acceptance} markers. \\
\path{docs/DEMO_GUIDE_week11_system_acceptance_testing.md} & New &
Section-by-section lecture companion for both readings. \\
\texttt{README.md} (suite and repository) & Modified & Layout, the
Week 11 entry, and the schedule. \\
\bottomrule
\end{tabular}
\end{center}

\begin{tipbox}
Install the new dependency before running anything this week:
\begin{lstlisting}[style=shellstyle]
cd demo-testing-suite
pip install -r requirements.txt
\end{lstlisting}
\end{tipbox}

% ============================================================
\section{The Readings' System, Runnable}
% ============================================================
\textbf{Location:} \path{demos/system_acceptance_testing/}. \textbf{Run
its commands from inside the folder:}
\begin{lstlisting}[style=shellstyle]
cd demos/system_acceptance_testing
pytest                         # 58 passed, in about ten seconds
\end{lstlisting}

\path{demos/integration_testing/} is unchanged: the Week 10 reading
quotes it line by line. The new folder started as a copy of it.

\begin{center}
\small
\begin{tabular}{@{}L{3.6cm}L{4.6cm}L{4.4cm}r@{}}
\toprule
\textbf{Folder} & \textbf{Level} & \textbf{How the shop runs} & \textbf{Tests}\\
\midrule
\path{tests/unit/} & unit (the developer's tests) & not at all & 8\\
\path{tests/integration/} & integration (Week 10's) & inside the test process & 13\\
\path{tests/system/} & system & separate process, HTTP only & 23\\
\path{tests/acceptance/} & acceptance (Gherkin) & separate process, HTTP only & 14\\
\bottomrule
\end{tabular}
\end{center}

\subsection*{Two defective components, one switch each}
\begin{lstlisting}[style=shellstyle]
SHOP_SERVICE=racy       pytest tests/system
SHOP_COUPONS=per_unit   pytest tests/acceptance
\end{lstlisting}
\begin{center}
\small
\begin{tabular}{@{}lcccc@{}}
\toprule
& \textbf{unit} & \textbf{integration} & \textbf{system} & \textbf{acceptance}\\
\midrule
\texttt{SHOP\_SERVICE=racy} & 8 passed & 13 passed & \textbf{1 failed} & 14 passed\\
\texttt{SHOP\_COUPONS=per\_unit} & 8 passed & 13 passed & 23 passed & \textbf{2 failed}\\
\bottomrule
\end{tabular}
\end{center}

Each defect is caught at exactly one level. The \emph{racy service}
charges the card before taking the stock, from a value it read before
the charge: with ten simultaneous customers it sells ten of five mugs
and leaves the stock at four. Only a test that runs the real,
multi-threaded server can see it. The \emph{per-unit pricing} reads
``FLAT5: \$5 off orders over \$10'' as ``\$5 off each unit over \$10''.
The developer's own tests order one unit, where both readings agree;
only the examples agreed with the product owner (two mugs; four
stickers) catch it.

% ============================================================
\section{A Real Finding: Hostile Input (F1)}
% ============================================================
When the system tests were first run against Week 10's input
validation, two hostile requests --- a \texttt{sku} sent as a JSON list,
and a \texttt{quantity} of $10^{30}$ --- crashed the request inside
SQLite, and Flask answered with its default \emph{HTML} \texttt{500}
page. That breaks the shop's own API contract (``every error response
is JSON\ldots never an HTML page''). It is present in
\path{demos/integration_testing/} as well; that folder is left as the
Week 10 reading printed it.

The fix needed a requirement that did not exist: the product owner set
a maximum of 100 units per order. \path{validate_order()} now checks
that \texttt{sku} is a string and \texttt{quantity} is 1--100, and a
JSON handler covers any \texttt{500}. To see the failure again:
\begin{lstlisting}[style=shellstyle]
patch -p1 < specs/f1_week10_validation.patch
pytest tests/system/test_robustness_and_security.py     # 2 failed
patch -R -p1 < specs/f1_week10_validation.patch
\end{lstlisting}

% ============================================================
\section{Updated: Expected Test Results}
% ============================================================
From the project root:
\begin{lstlisting}[style=shellstyle]
cd demo-testing-suite
pip install -r requirements.txt
pytest -v
\end{lstlisting}
you should still see \textbf{75 passed, 3 xfailed, 1 failed}, exactly
Week 10's result: nothing in \path{app/} or \path{tests/} changed this
week. \path{demos/system_acceptance_testing/} is not part of that run
(\texttt{pytest.ini} sets \texttt{testpaths = tests}); run it from inside
its folder.

\begin{warningbox}
\textbf{Why the main suite gets no system tests this week.}
\path{app/} is a library of modules, not a program that runs on its
own: there is no system to start. System and acceptance testing need a
deployable system, which is why both readings use the shop. Your
semester project, which is a runnable system, is where you apply them.
\end{warningbox}

% ============================================================
\section{How This Connects to the Readings}
% ============================================================
\begin{keyidea}
Weeks 2--10 tested the system's parts and the promises between them.
Week 11 tests the whole: first against the team's specification, as a
customer, operator, or attacker would meet it (system testing); then
against what the customer actually meant (acceptance testing). Each
level found a defect every level below it passed.
\end{keyidea}

% ============================================================
\section{Suggested Next Steps}
% ============================================================
\begin{enumerate}
\item Run the switch table above yourself, one cell at a time, and
\emph{predict} each result first.
\item Read \path{shop/service_racy.py} and find the defect before running
the system tests against it.
\item Before reading \path{specs/user_stories.md}'s three-amigos notes,
write down what you think 2 mugs with FLAT5 should cost.
\item Implement restocking (System Testing, Exercise 6) and order
history (Acceptance Testing, Exercise 6), tests first.
\item Write a system test plan and three Gherkin scenarios for your
semester project.
\end{enumerate}

\end{document}
